Large language model (LLM)–based agent applications often incur high response time. Speculative decoding is a promising solution to improve the inference efficiency of LLM agents without impacting generation quality. 
%
However, existing speculative decoding algorithms exhibit substantial speed degradation as batch sizes grow, limiting their practicality to deploy in real-world agent applications.
%
In this work, we first present a systematic analysis of speculative decoding for LLM agents and identify two dominant factors of speedup degradation: high rejection rate of speculative tokens, and under-utilization of dynamic token budgets.
%
Motivated by these findings, we propose \sysname, a speculative decoding algorithm that addresses the limitations of existing methods for LLM agents. \sysname incorporates structure-isolated drafting that constrains speculation to semantically coherent segments of the agent workflow, reducing the drafts of irrelevant semantic paths and achieving an extremely low rejection rate. Moreover, \sysname adopts redundancy-aware budget allocation that exploits agent-level information to better utilize the dynamically-free token budget during the agent inference.
%
We implement and evaluate \sysname on five different workloads and four different models from four different LLM families in vLLM. Our results demonstrate the superiority of \sysname over state-of-the-art.
%
% The source code is available at \url{https://github.com/AIoT-MLSys-Lab/SVD-LLM}.
